В данном разделе доказывается, что величина, вычисляемая методом оконного ДП из раздела~2, является строгой нижней оценкой константы $\gamma_\sigma$, и анализируется вычислительная сложность алгоритма по памяти и времени.

\smallskip

\textbf{Обозначения.} Через $\gamma_\sigma$ обозначим истинную константу Чватала--Санкоффа для алфавита $\Sigma_\sigma$ (см.\,определение~1.1, \S1.3). Через $\gamma_{\sigma, N, \mathrm{MAX\_SH}}$ --- предел нормированной ДП-функции при фиксированных параметрах окна $N$ и максимального сдвига $\mathrm{MAX\_SH}$ (определяется ниже в теореме~3.4). Для произвольной допустимой стратегии $\pi$ обозначим $\gamma^\pi := \lim_{T \to \infty} \mathbb{E}[S_T^\pi]/T$, если предел существует.

\subsection{Лемма о нижней оценке через стратегию}

\textbf{Лемма 3.1.} \emph{Пусть $\pi$ --- произвольная допустимая стратегия с ограниченным обзором (определена в~\S2.1), $T$ --- число её шагов, $S_T^\pi$ --- длина построенной общей подпоследовательности, $T_A, T_B$ --- финальные значения указателей. Тогда}
\[
    S_T^\pi \leq \mathrm{lcs}\bigl(A_{1..T_A},\, B_{1..T_B}\bigr).
\]

\textbf{Доказательство.} По построению стратегия $\pi$ строит конкретную последовательность совпадений символов $a_{i_1} = b_{j_1}, a_{i_2} = b_{j_2}, \ldots, a_{i_k} = b_{j_k}$, где $i_1 < i_2 < \ldots < i_k$ и $j_1 < j_2 < \ldots < j_k$: действия~(L), (R) и засчитывание совпадения строго увеличивают соответствующий указатель и никогда не уменьшают его, поэтому индексы совпадений строго монотонно возрастают по обеим строкам. Эта последовательность задаёт общую подпоследовательность строк $A_{1..T_A}$ и $B_{1..T_B}$ длины $S_T^\pi = k$. Так как $\mathrm{lcs}$ --- максимум по всем общим подпоследовательностям, требуемое неравенство выполнено. \hfill$\square$

\subsection{Связь со средним}

Применяя лемму~3.1 к случайным строкам $A, B$ и беря математическое ожидание по обеим строкам, получаем:
\[
    \mathbb{E}_\pi[S_T^\pi] \leq \mathbb{E}\bigl[\mathrm{lcs}(A_{1..T_A},\, B_{1..T_B})\bigr] \leq \mathbb{E}\bigl[\mathrm{lcs}(A_{1..T},\, B_{1..T})\bigr] = \mathbb{E}[L_T].
\]
Второе неравенство --- монотонность LCS по префиксу: для любых строк $X, X', Y, Y'$ с $X \preceq X'$, $Y \preceq Y'$ (где $\preceq$ обозначает префиксное вложение) любая общая подпоследовательность пары $(X, Y)$ является общей и для пары $(X', Y')$, откуда $\mathrm{lcs}(X, Y) \leq \mathrm{lcs}(X', Y')$. При $T_A \leq T$, $T_B \leq T$ это и даёт требуемое. Деление на $T$ и переход к пределу $T \to \infty$ даёт
\[
    \limsup_{T \to \infty} \frac{\mathbb{E}[S_T^\pi]}{T} \leq \gamma_\sigma.
\]

\textbf{Следствие 3.2.} \emph{Любая стратегия $\pi$ с ограниченным обзором, дающая $\mathbb{E}[S_T^\pi] / T \to \gamma^\pi$ при $T \to \infty$, удовлетворяет $\gamma^\pi \leq \gamma_\sigma$.}

\subsection{Теорема о корректности ДП}

\textbf{Теорема 3.3} (корректность ДП).
\emph{Пусть $f_K(\mathrm{sh}, a, b)$ определена рекуррентным соотношением из~\S2.3 с граничными условиями~\S2.4. Тогда $f_K(\mathrm{sh}, a, b)$ есть в точности максимальное по всем допустимым стратегиям $\pi$ значение $\mathbb{E}_\pi[S_K^\pi]$, при условии что в начальный момент состояние есть $(\mathrm{sh}, a, b)$, и за $K$ шагов в концы окон поступают независимые равномерные на $\Sigma_\sigma$ символы.}

\textbf{Доказательство.} Индукция по $K$.

\emph{База ($K = 0$):} стратегия не делает ни одного шага, $S_0^\pi = 0$. Значение $f_0 = 0$ совпадает.

\emph{Шаг.} Предположим, утверждение верно для $K-1$, докажем для $K$.
Рассмотрим состояние $(\mathrm{sh}, a, b)$. Пусть $(a_0, b_0)$ --- новые символы, поступающие в окна на текущем шаге, $a_0, b_0 \sim U(\Sigma_\sigma)$ независимо. Обозначим младшие символы окон $u = a \bmod \sigma$ (это $a_i$), $v = b \bmod \sigma$ (это $b_j$). Возможны два случая.

\textit{Случай 1: $u = v$.} По правилам стратегии (см.~\S2.1) алгоритм обязан засчитать совпадение пары $(u, v)$ и сдвинуть оба указателя. Новое состояние --- $(\mathrm{sh},\, a',\, b')$, где $a' = (a \gg \mathrm{BITS}) \,|\, (a_0 \ll (\mathrm{MASK\_BITS} - \mathrm{BITS}))$, аналогично $b'$. Текущий вклад --- $+1$. По индукционному предположению ожидаемое продолжение есть $f_{K-1}(\mathrm{sh}, a', b')$. Полное значение:
\[
    T_K = 1 + f_{K-1}(\mathrm{sh}, a', b'),
\]
что совпадает с формулой~\S2.3.

\textit{Случай 2: $u \neq v$.} Стратегия выбирает между действиями (L) и (R). Действие (L) переводит в состояние $(\mathrm{sh}-1, a', b)$, действие (R) --- в $(\mathrm{sh}+1, a, b')$. Текущий вклад --- $0$. Оптимальный выбор:
\[
    T_K = \max\bigl(f_{K-1}(\mathrm{sh}-1, a', b),\, f_{K-1}(\mathrm{sh}+1, a, b')\bigr).
\]
При $\mathrm{sh} = 0$ или $\mathrm{sh} = \mathrm{MAX\_SH}$ соответствующая ветвь недоступна и максимум вырождается в одно слагаемое, как описано в~\S2.4.

Так как $(a_0, b_0)$ независимы от прошлого, операция $\mathbb{E}[\,\cdot\,]$ коммутирует с детерминированным $\max$ при условии независимости. Полное:
\[
    f_K(\mathrm{sh}, a, b) = \mathbb{E}_{a_0, b_0}[T_K] = \sum_{a_0, b_0 \in \Sigma_\sigma} \frac{1}{\sigma^2} \cdot T_K,
\]
что и доказывает шаг индукции. \hfill$\square$

\subsection{Теорема о сходимости}

\textbf{Теорема 3.4} (нижняя оценка).
\emph{Обозначим $s_0 := (0, 0, 0)$ --- начальное состояние ДП. Тогда:}

\emph{(i) для любого $K \geq 1$ выполнено $f_K(s_0) / K \leq \mathbb{E}[L_K]/K$;}

\emph{(ii) в частности,
\[
    \limsup_{K \to \infty} \frac{f_K(s_0)}{K} \;\leq\; \gamma_\sigma;
\]
то есть для любого $\varepsilon > 0$ и достаточно большого $K$ выполнено $f_K(s_0)/K \leq \gamma_\sigma + \varepsilon$.}

\textbf{Доказательство.} По теореме~3.3 значение $f_K(s_0)$ равно $\mathbb{E}_\pi[S_K^\pi]$ для оптимальной стратегии $\pi = \pi^\ast_K$ с горизонтом $K$, стартующей из состояния $s_0$ на случайных строках $A, B$. Эта стратегия допустима в смысле определения~\S2.1, поэтому к ней применима лемма~3.1 и вывод из~\S3.2 о связи со средним:
\[
    \frac{f_K(s_0)}{K} \;=\; \frac{\mathbb{E}[S_K^{\pi^\ast_K}]}{K} \;\leq\; \frac{\mathbb{E}[L_K]}{K}.
\]
Это даёт (i). Утверждение (ii) получается переходом к пределу $K \to \infty$ и применением теоремы~1.1: $\mathbb{E}[L_K]/K \to \gamma_\sigma$. \hfill$\square$

\smallskip

\textbf{Замечание.} Существование самого предела $\lim_{K \to \infty} f_K(s_0)/K$ для обоснования нижней оценки $\gamma_\sigma$ не требуется: согласно пункту~(i) теоремы~3.4, \emph{каждое} конкретное значение $f_K(s_0)/K$, полученное при фиксированном $K$, уже не превосходит $\mathbb{E}[L_K]/K \leq \gamma_\sigma$ (с точностью до известной эмпирической оценки разности $\gamma_\sigma - \mathbb{E}[L_K]/K$, см.~\S5.4). Поэтому каждое отдельное вычисление $f_K(s_0)/K$ даёт строгую нижнюю оценку константы Чватала--Санкоффа независимо от вопроса о существовании предела.

\textbf{Следствие 3.5.} \emph{Для любых $N \geq 1$, $\mathrm{MAX\_SH} \geq 0$ и достаточно большого $K$ выполнено}
\[
    \frac{f_K(0, 0, 0)}{K} \leq \gamma_\sigma + \varepsilon
\]
с погрешностью $\varepsilon = \varepsilon(K) \to 0$. На практике сходимость к пределу наблюдается экспоненциально быстро (см.\,раздел~5).

\subsection{Анализ вычислительной сложности}

\textbf{Память.}
Для слоя ДП фиксированного значения $\mathrm{sh}$ требуется массив размера $\sigma^N \times \sigma^N = M^2$, где $M = \sigma^N$. Хранятся слои для всех $\mathrm{sh} \in [0, \mathrm{MAX\_SH}]$ --- это $\mathrm{MAX\_SH} + 1$ слой --- и два вспомогательных «зеркальных» слоя $m_0, m_1$ для корректной обработки граничных случаев $\mathrm{sh} = 0$ и $\mathrm{sh} = \mathrm{MAX\_SH}$ без условных переходов (см.~\S2.4). Итого хранится $\mathrm{MAX\_SH} + 1 + 2 = \mathrm{MAX\_SH} + 3$ слоёв; в асимптотике это $O(\mathrm{MAX\_SH} \cdot \sigma^{2N})$ элементов. Запись:
\[
    \text{память} = O\bigl((\mathrm{MAX\_SH} + 2) \cdot \sigma^{2N}\bigr) \text{ элементов float} = O\bigl(\mathrm{MAX\_SH} \cdot \sigma^{2N}\bigr) \text{ байт} \cdot 4.
\]
Для $\sigma = 2$, $N = 12$, $\mathrm{MAX\_SH} = 30$: $\sigma^{2N} = 2^{24} \approx 1{,}68 \cdot 10^7$ элементов на слой; всего $31 \cdot 2^{24} \cdot 4 \approx 2$~ГБ.

\textbf{Время одной итерации ДП.}
На каждом шаге $K \to K+1$ необходимо обновить все состояния. Обновление одного состояния $(\mathrm{sh}, a, b)$ требует $O(\sigma^2)$ обращений к ячейкам предыдущего слоя (сумма по $(a_0, b_0) \in \Sigma_\sigma^2$). Итого:
\[
    \text{время итерации} = O\bigl(\mathrm{MAX\_SH} \cdot \sigma^{2N + 2}\bigr).
\]

\textbf{Полное время.}
\[
    \text{полное время} = O\bigl(K \cdot \mathrm{MAX\_SH} \cdot \sigma^{2N + 2}\bigr).
\]

\textbf{Иллюстрация.}
Для $\sigma = 2$, $N = 12$, $\mathrm{MAX\_SH} = 30$, $K = 40\,000$: грубая оценка $4 \cdot 10^4 \cdot 31 \cdot 4 \cdot 2^{24} \approx 8 \cdot 10^{13}$ элементарных операций. С учётом эффективного использования AVX2 (8 одинарных операций за такт), кэш-локальности и развёртки циклов это реализуется за $\approx 5$ часов на одном ядре современного процессора.

\subsection{Источники численной погрешности}

При реализации используется одинарная точность (\texttt{float}), что даёт около $7$ знаков точности. Контрольный пересчёт всей задачи в \texttt{double} (на меньших параметрах, где это осуществимо) подтверждает совпадение $5$--$6$ значащих знаков с \texttt{float}-вариантом. Округление в третий-четвёртый знак после запятой не влияет на формулировку результата $\gamma_2 \geq 0{,}7964$, поскольку зазор до точного предела с запасом превышает погрешность округления.

\subsection*{Выводы по разделу}
\addcontentsline{toc}{subsection}{Выводы по разделу 3}

\begin{enumerate}
    \item Доказана лемма~3.1, согласно которой длина общей подпоследовательности, построенной стратегией, не превосходит LCS.
    \item Доказана теорема~3.3 о корректности ДП-уравнений (индукция по числу шагов).
    \item Доказана теорема~3.4 о нижней оценке: для любого $K \geq 1$ выполнено $f_K(0,0,0)/K \leq \mathbb{E}[L_K]/K$, откуда $\limsup_K f_K(0,0,0)/K \leq \gamma_\sigma$. Это означает, что вычисляемая методом величина является строгой нижней оценкой $\gamma_\sigma$ при \emph{каждом} фиксированном $K$.
    \item Память алгоритма $O(\mathrm{MAX\_SH} \cdot \sigma^{2N})$, время $O(K \cdot \mathrm{MAX\_SH} \cdot \sigma^{2N+2})$. На рекордных параметрах $\sigma = 2$, $N = 12$, $\mathrm{MAX\_SH} = 30$ это $\approx 2$~ГБ и часы счёта.
\end{enumerate}
